\documentclass{article}
\usepackage{geometry}
\geometry{margin=1.5cm, vmargin={0pt,1cm}}
\setlength{\topmargin}{-1cm}
\setlength{\headheight}{12.5pt}
\setlength{\paperheight}{29.7cm}
\setlength{\textheight}{25.3cm}

% useful packages.
\usepackage[UTF8]{ctex}
\usepackage{amsfonts}
\usepackage{amsmath}
\usepackage{amssymb}
\usepackage{amsthm}
\usepackage{enumerate}
\usepackage{graphicx}
\usepackage{multicol}
\usepackage{fancyhdr}
\usepackage{layout}
\usepackage{luacode}
\usepackage[ruled,linesnumbered]{algorithm2e}

\newcommand{\getEnv}[2]{\directlua{tex.print(os.getenv("#1") or "#2")}}

% some common command
\newcommand{\dif}{\mathrm{d}}
\newcommand{\innerProd}[1]{\left\langle #1 \right\rangle}
\newcommand{\difFrac}[2]{\frac{\dif #1}{\dif #2}}
\newcommand{\pdfFrac}[2]{\frac{\partial #1}{\partial #2}}
\newcommand{\T}{\mathrm{T}}

% Name and uID
\newcommand{\name}{\getEnv{NAME}{NAME}}
\newcommand{\uid}{\getEnv{UID}{UID}}
\newcommand{\course}{数值代数}
\newcommand{\hwNum}{1}

\newcommand{\vecTwo}[2]{
  \begin{bmatrix}
    #1 \\
    #2
\end{bmatrix}}
\newcommand{\vecThree}[3]{
  \begin{bmatrix}
    #1 \\
    #2 \\
    #3
\end{bmatrix}}
\newcommand{\vecTwoH}[2]{
  \begin{bmatrix}
    #1 & #2
  \end{bmatrix}
}
\newcommand{\vecThreeH}[3]{
  \begin{bmatrix}
    #1 & #2 & #3
  \end{bmatrix}
}

\newenvironment{solution}{
\begin{proof}[解]}{
\end{proof}}

\begin{document}

\pagestyle{fancy}
\fancyhead{}
\lhead{\name (\uid)}
\chead{\course \#\hwNum}
\rhead{\today}

\section{习题 1.1}

给出求下三角阵的逆的详细算法

\begin{solution}
  \mbox{} \\
  \begin{algorithm}[H]
    \caption{下三角矩阵求逆}
    \KwIn{下三角矩阵 $L \in \mathbb{R}^{n \times n}$}
    \KwOut{$L^{-1}$}

    \For{$j = 1 : n$}{
      $L^{-1}(j,j) = \frac{1}{L(j,j)}$\;

      \For{$i = j+1 : n$}{
        $s = L(i,:) \cdot L^{-1}(:,j)$\;
        $L^{-1}(i,j) = -\frac{s}{L(i,i)}$\;
      }
    }

    \Return{$L^{-1}$}\;
  \end{algorithm}
\end{solution}

\section{习题 1.4}
确定 $3 \times 3$ Gauss 变换 $L$ 使得
\[
  L \vecThree{2}{3}{4} = \vecThree{2}{7}{8}
\]

\begin{solution}

  $L = I - l_{k} e_k^T$ ，从而有

  \[
    \begin{aligned}
      L \vecThree{2}{3}{4}
      &= \vecThree{2}{3}{4} - l_k e_k^T \vecThree{2}{3}{4} \\
      &= \vecThree{2}{7}{8}
    \end{aligned}
  \]

  即：

  \[
    l_k e_k^T \vecThree{2}{3}{4} = \vecThree{2}{3}{4} - \vecThree{2}{7}{8} = \vecThree{0}{-4}{-4}
  \]

  从而得到 $k = 1$ ，于是有 $2l_k = (0,-4,-4)^T$ ，从而有 $l_k = (0,-2,-2)^T$ ，从而：

  \[
    L =
    \begin{bmatrix}
      1 & 0 & 0 \\
      2 & 1 & 0 \\
      2 & 0 & 1
    \end{bmatrix}
  \]

\end{solution}

\section{习题 1.7}
设 $A$ 对称且 $a_{11}\neq 0$，并假定经过一步 Gauss 消去之后，$A$ 具有如下形状：
\[
  \begin{bmatrix}
    a_{11} & a_1^{\T} \\
    0 & A_2
  \end{bmatrix}
\]
证明：$A_2$ 仍是对称矩阵。

\begin{proof}
  设第一步 Gauss 消去矩阵为 $L_1$，于是有
  \[
    L_1 A = (I - l_1 e_1^{\T}) A = A - l_1
    \begin{bmatrix}
      a_{11} & a_{1}^{\T}
    \end{bmatrix} =
    \begin{bmatrix}
      a_{11} & a_1^{\T} \\
      0 & A_2
    \end{bmatrix}
  \]
  又有：
  \[
    l_1 = \vecTwo{0}{a_1 / a_{11}}
  \]
  于是有：
  \[
    l_1 \vecTwoH{a_{11}}{a_1^{\T}}
    = \vecTwo{0}{a_1 / a_{11}} \vecTwoH{a_{11}}{a_1^{\T}} =
    \begin{bmatrix}
      0 & 0 \\
      a_1 & a_1 a_1^{\T} / a_{11}
    \end{bmatrix}
  \]
  其中 $a_1 a_1^{\T} / a_{11}$ 显然为对称矩阵，且原矩阵 $A$ 的子矩阵 $A[2:,2:]$ 同样为对称矩阵，于是有
  \[
    A_2 = A[2:,2:] - a_1 a_1^{\T}
  \]
  也是对称矩阵。

\end{proof}

\section{习题 1.8}
设 $A=[a_{ij}]\in\mathbb{R}^{n\times n}$ 是严格对角占优阵，即对每个 $k=1,\dots,n$ 有
\[
  |a_{kk}|>\sum_{j=1,j\neq k}^n |a_{kj}|.
\]
又设经过一步 Gauss 消去之后，$A$ 具有如下形状：
\[
  \begin{bmatrix}
    a_{11} & a_1^{\mathrm{T}} \\
    0 & A_2
  \end{bmatrix}
\]
试证：矩阵 $A_2$ 仍是严格对角占优阵。

\begin{proof}
  设原矩阵为：

  \[
    A =
    \begin{bmatrix}
      a_{11} & a_1^{\T} \\
      a'_1 & A_1
    \end{bmatrix}
  \]
  于是有：

  \[
    L_1 A = (I - l_1 e_1^{\T}) A
    = A - \vecTwo{0}{a'_1/a_{11}} \vecTwoH{a_{11}}{a_1^{\T}}
  \]

  从而：

  \[
    A_2 = A_1 - \frac{1}{a_{11}} a'_1 a_1^{\T}
  \]

  记：
  \[
    B =
    \begin{bmatrix}
      1 & 0^{\T} \\
      0 & A_2 \\
    \end{bmatrix}
  \]
  于是有：
  \[
    b_{ij} =
    \begin{cases}
      1 & i = j = 1 \\
      0 & i = 1 \ \text{or}\ j = 1 \\
      a_{ij} - a_{i1} a_{1j} / a_{11} & \text{otherwise}
    \end{cases}
  \]
  只需证明 $B$ 严格对角占优即可。对于第一行显然，于是只需证明：

  \[
    |b_{kk}| = |a_{kk} - a_{k1}a_{1k} / a_{11}| > \sum_{j=2, j\neq k}^{n} |a_{kj} - a_{k1}a_{1j} / a_{11}|
  \]

  由于 $A$ 严格对角占优，故有：

  \[
    \begin{aligned}
      |a_{kk} - a_{k1}a_{1k} / a_{11}|
      &\geq |a_{kk}| - |a_{k1}a_{1k} / a_{11}| \\
      &> \sum_{j=1, j\neq k}^{n} |a_{kj}| - |a_{k1}||a_{1k}| / |a_{11}| \\
      &= \sum_{j=2, j\neq k}^{n} |a_{kj}| + |a_{k1}| (1 - \frac{|a_{1k}|}{|a_{11}|}) \\
      &= \sum_{j=2, j\neq k}^{n} |a_{kj}| + \frac{|a_{k1}|}{|a_{11}|} (|a_{11}| - |a_{1k}|) \\
      &> \sum_{j=2, j\neq k}^{n} |a_{kj}| + \frac{|a_{k1}|}{|a_{11}|} (\sum_{j=2, j\neq k}^{n} |a_{1j}|) \\
      &= \sum_{j=2, j\neq k}^{n} (|a_{kj}| + \frac{|a_{k1}a_{1j}|}{|a_{11}|}) \\
      &\geq \sum_{j=2, j\neq k}^{n} |a_{kj} - a_{k1}a_{1j} / a_{11}|
    \end{aligned}
  \]

\end{proof}

\section{习题 1.10}
设 $A$ 是正定阵。如果对 $A$ 执行 Gauss 消去一步产生一个形式为
\[
  \begin{bmatrix}
    a_{11} & a_1^{\mathrm{T}} \\
    0 & A_2
  \end{bmatrix}
\]
的矩阵，证明：$A_2$ 仍是正定阵。

\begin{proof}
  设
  \[
    A =
    \begin{bmatrix}
      a_{11} & r^{\T} \\
      r & B
    \end{bmatrix}
  \]
  其中 $r = a_1$ 。为了避免混淆，后面都用 $r$ 而非 $a_1$。于是有：
  \[
    A_2 = B - \frac{1}{a_{11}}r r^{\T}
  \]

  考虑二次型 $f(y) = y^{\T} A_2 y^{\T}$ ，于是有：

  \[
    f(y) = y^{\T} (B - \frac{1}{a_{11}} r r^{\T}) y = y^{\T}By - \frac{1}{a_{11}} (r^{\T}y)^2
  \]

  令 $x = \vecTwo{x_1}{y}$ ，于是有：

  \[
    \begin{aligned}
      x^{\T}Ax &= \vecTwoH{x_1}{y^{\T}}
      \begin{bmatrix}
        a_{11} & r^{\T} \\
        r & B
      \end{bmatrix} \vecTwo{x_1}{y} \\
      &= a_{11} x_1^2 + 2x_1 r^{\T} y + y^T B y \\
      &= a_{11} (x_1 + \frac{r^{\T}y}{a_{11}})^2 + y^{\T} B y - \frac{1}{a_{11}} (r^{\T}y)^2 \\
      &\geq y^{\T} B y - \frac{1}{a_{11}} (r^{\T}y)^2 = f(y)
    \end{aligned}
  \]

  由于 $A$ 正定，故 $x^{\T} A x > 0$ 对任意非零 $x$ 成立，于是 $f(y) > 0$ 对任意非零 $y$ 成立，即 $A_2$ 正定。

\end{proof}

\end{document}
